\providecommand{\PushA}{\frac{5422587}{6103515625}}
\providecommand{\PushK}{890666}
\providecommand{\PushTheta}{0.99999999997327}
\section{From the five estimates to Theorem~C}
\label{sec:comparison-proof}

\paragraph{What this section proves and how to read it.}
This section proves Theorem~C. The closing inequality
(\cref{thm:record:criterion}) combines the wrapping bound $W$, the mean energies $e_{\rm L},e_{\rm R}$ and the
envelope energies $V_{\rm L},V_{\rm R}$ of the previous section with the first comparison: at a
least counterexample, $1\le\Theta+o(1)$. The robustness lemma
(\cref{lem:record:rec}) shows that every $o(1)$ in this chain is produced by
structured energies and by the two analytic remainders, so that the
density-linear induction of \cref{thm:record:lineardensity} applies. The
section ends with the proof of the one-query proposition
(\cref{rem:cp:one-query}) and with \cref{tab:hypotheses}, which records where
Part~V verifies each hypothesis.

\subsection{At a least counterexample, \texorpdfstring{$1\le\Theta+o(1)$}{1 <= Theta+o(1)}}\label{subsec:cp:closing}

This lemma is where the comparison closes: the five estimates of
\cref{sec:comparison-lemmas} combine to show that a least
counterexample forces $1\le\Theta+o(1)$. Hence $\Theta<1$ rules out
a counterexample at large scale, which is the contradiction behind
Theorem~C.

\begin{lemma}[Closing inequality]\label{thm:record:criterion}
Assume (CP0)--(CP4), (CP6) and (CP8). Let $\mathcal E$ and $m$ be numbers such
that, at every least normalized counterexample,
\begin{enumerate}[label=(\roman*)]
\item $Y_\nu\ge Y_\lambda-(\mathcal E+o(1))Y$, and
\item $\frac12(e_{\rm L}+e_{\rm R})\ge(m-o(1))Y$.
\end{enumerate}
Then, at every least normalized counterexample,
\begin{equation}\label{eq:record:criterion}
 1\le\Theta(\mathcal E,m)+o(1),\qquad
 \Theta(\mathcal E,m)=\bigl(1-2^{-k}\bigr)\bigl(L_0+\mathcal E\bigr)+J+W
 +\frac{1-m}{2^k},
\end{equation}
with $J$ and $W$ as in \cref{thm:comparison}. Consequently, if
$\Theta(\mathcal E,m)<1$, no least normalized counterexample
exists with $\Lsc$ sufficiently large.
\end{lemma}

\begin{proof}
By \cref{lem:record:wrap-expansion},
\begin{equation}\label{eq:cp:wrap-step}
 2^kY_\nu\le J_\nu+(2^kW+o(1))Y .
\end{equation}
The index-support lemma, in the form of \cref{cor:index-support-frame}, applies
to the signed index operator $\mathcal T=\bigotimes_e(2P_e-T_e)$, the joint
bracket of the envelopes of (CP6) tensorized over the finitely many atoms of
$\mathcal A_N$, and the half-functions $X_\omega$, $Y_\omega$, whose means are
$\bar X$. Its mean term is $Y_\nu$, its mean energies are the $e_s$ of (CP7),
and its random-inclusion energies are $V_{\rm L}=V_{\bigotimes M_{\rm L}}(X)$ and
$V_{\rm R}=V_{\bigotimes M_{\rm R}}(Y)$ (\cref{def:hyb:coupled}(f)). By
\eqref{eq:index-support-frame},
\begin{equation}\label{eq:cp:mean-step}
 J_\nu\le Y_\nu+\tfrac12\bigl[(V_{\rm L}-e_{\rm L})+(V_{\rm R}-e_{\rm R})\bigr].
\end{equation}
By \cref{lem:envelope-energy}, applied with $Q_{\Blk}=M_{s,\Blk}-\Nbr_{\Blk}$
and $T_Q=T$ (CP6), $V_{\rm L},V_{\rm R}\le(Be^TR_I^k+o(1))Y=(2^kJ+o(1))Y$.
With (ii), \eqref{eq:cp:wrap-step} and \eqref{eq:cp:mean-step} give
\[
 (2^k-1)\,Y_\nu\le\bigl(2^kJ-m+2^kW+o(1)\bigr)Y .
\]
By (i) and \cref{lem:record:first-stage}, $Y_\nu\ge(1-L_0-\mathcal E-o(1))Y$.
Since $Y>0$ and $2^k-1>0$,
$(2^k-1)(1-L_0-\mathcal E)\le2^kJ-m+2^kW+o(1)$, that is,
$2^k\le(2^k-1)(L_0+\mathcal E)+2^kJ+2^kW+(1-m)+o(1)$; dividing by $2^k$ gives
\eqref{eq:record:criterion}. The last assertion follows because $o(1)\to0$ as
$\Lsc\to\infty$, uniformly over square-difference-free $A\subseteq[\Kint\Lsc]$.
\end{proof}

\begin{corollary}[Joint comparison and mean allocation]\label{cor:joint-closing}
The comparison loss $Y_\lambda-Y_\nu$ and the mean energies
$e_{\rm L},e_{\rm R}$ enter the preceding proof with opposite signs.
It suffices to bound their combined contribution.
The separate bounds (i) and (ii) in \cref{thm:record:criterion} may be
replaced by a single robust bound. Suppose that, for some real
constant $\mathcal H$,
\[
 (2^k-1)(Y_\lambda-Y_\nu)-\tfrac12(e_{\rm L}+e_{\rm R})
 \le (2^k\mathcal H-1+o(1))Y.
\]
Then the resulting closing number is
\[
 \Theta=(1-2^{-k})L_0+J+W+\mathcal H.
\]
In particular, positive comparison energies and negative mean energies
may first be combined on the same labelled diagonal. If their residual
coefficients are nonnegative, their common energies may then be bounded
by upper prices and completed to exponential series.
\end{corollary}

\begin{proof}
Keep the mean energies in \eqref{eq:cp:mean-step} until after combining
that inequality with \eqref{eq:cp:wrap-step}. This gives
\[
 (2^k-1)Y_\lambda\le (2^k(J+W)+o(1))Y
 +(2^k-1)(Y_\lambda-Y_\nu)-\tfrac12(e_{\rm L}+e_{\rm R}).
\]
Use the displayed joint bound and
$Y_\lambda\ge(1-L_0-o(1))Y$, and divide by $2^k$.
Nonnegative residual coefficients permit exactly the positive
completion argument of \cref{lem:record:record-expansion}; applying an
upper price to an isolated negative mean term would not permit it.
The robustness proof of \cref{lem:record:rec} applies to the combined
positive series, with the same two kinds of analytic remainder.
\end{proof}

\begin{remark}[The trivial value of the mean constant]\label{rem:cp:trivial-mean}
By \cref{lem:record:domination}, hypothesis (ii) always holds with
$m=1-L_0$, and then $\Theta(\mathcal E,1-L_0)=L_0+(1-2^{-k})\mathcal E+J+W$;
that remark explains this bookkeeping. A larger $m$ lowers $\Theta$ by
$(m-1+L_0)/2^k$; Part~V proves $m>1$ (\cref{prop:record:mean-credit}).
\end{remark}

\subsection{Every error term comes from the two allowed remainders}\label{subsec:cp:robust}

The hypothesis (Rec) of \cref{thm:record:lineardensity} asks for the closing
inequality with every analytic error made explicit: every structured-energy
bound inflated by the factor $1+z$, every remainder bounded by $zY$, and a
closing number $\Theta(z)$ continuous at $z=0$. The next lemma supplies it and
records how each $o(1)$ above arises.

\begin{lemma}[The recurrence hypothesis]\label{lem:record:rec}
Assume (CP0)--(CP8).
\begin{enumerate}[label=(\roman*)]
\item\emph{Robust bounds.} Exact same-scale relations (identities and
 inequalities between quantities at the same scale, with no error term) hold
 unchanged at every $z$-inflated counterexample. A bound $\Xi\le(c+o(1))Y$ obtained by applying \cref{prop:record:density-structured}, directly or
 through \cref{lem:record:folded-energy}, \ref{it:record:E1} or
 \cref{lem:bracket-energy}, to finitely many families of structured
 energies, summing with nonnegative coefficients bounded by convergent
 products of prices, and adding remainders of the forms of
 \cref{def:robust}(a), is robust: at a $z$-inflated counterexample its
 constant $c$ becomes $(1+z)c+c'z$, where $c'$ is a bound, independent of $\Lsc$ and $A$, for the coefficients multiplying the remainder integrals and sums in its derivation. An exact relation minus robust upper bounds is a robust lower
 bound. Let a bound be derived from finitely many robust bounds, with
 functions $c_1(z),\dots,c_n(z)$, by an argument that is valid whenever
 finitely many strict inequalities between the $c_i(z)$ and fixed numbers
 hold, and let its constant be $F(c_1,\dots,c_n)$, with $F$ continuous at
 $(c_1(0),\dots,c_n(0))$. If the strict inequalities hold at $z=0$, the bound
 is robust, with $c(z)=F(c_1(z),\dots,c_n(z))$.
\item\emph{The lemmas.} The bounds of \cref{lem:energy-toolkit} and of
 \cref{lem:sharp-first,lem:record:first-stage,lem:bracket-energy,lem:record:wrap-expansion,lem:record:common-mean,cor:index-support-frame,lem:record:record-expansion,lem:record:variance-interface,lem:envelope-energy}
 are exact or robust for fixed selected sets. For the summed families used
 here, the full-union split of \cref{sec:comparison-lemmas} gives uniform
 $z_0$, $\Lsc_1$ and aggregate remainder coefficients. Their remainders
 are of two kinds:
 \begin{enumerate}[label=(\alph*)]
 \item the nonwrapping terms of \cref{lem:record:wrap-expansion}: for each of
  the at most $3^kn_\nu$ triples of a product component of $\nu$, an index
  subset and a subset of wrapping queries that does not contain every query of the subset (a \emph{wrap pattern with a nonwrapping query}), an integral of the form
  bounded in \cref{lem:record:density-can}, against a probability law of the
  indices, the coefficients of the product components for one pair $(X_\omega,Y_\omega)$ of half-functions adding up to at most $2^k\mathfrak v_\nu e^{P+S}$;
 \item the terms of the selections whose union $E$ of atoms has $q_E>N^\tau$,
  in finitely many sums of the form bounded in \cref{lem:record:density-lp}
  whose coefficients $\eta_{\Blk}$ are the prices of the selected laws after the operator-to-energy reduction of \cref{lem:record:rec}(ii) (the \emph{post-reduction prices}); together they have the bound of one such sum, with $c$ the sum of their constants, and with $\Pi_{\rm H}=\prod_\Blk(1+\eta_\Blk q_\Blk^{1/4})$ as in \cref{lem:record:density-lp}; neither $c$ nor $\Pi_{\rm H}$ depends on $\Lsc$ or $A$.
 \end{enumerate}
\item\emph{The recurrence.} Let $\mathcal E$ and $m$ be numbers for which (i)
 and (ii) of \cref{thm:record:criterion} hold robustly. Then there are
 $z_0>0$, $\Lsc_1$ and a function $\Theta\colon[0,z_0]\to[0,\infty)$,
 continuous at $0$ with $\Theta(0)=\max(\Theta(\mathcal E,m),0)$, such that
 the hypothesis (Rec) of \cref{thm:record:lineardensity} holds; its
 remainders are those of (ii) together with those of the robust arguments
 for (i) and (ii) of \cref{thm:record:criterion}, and they have the forms
 that theorem requires.
\end{enumerate}
\end{lemma}

\begin{proof}
(i) An exact relation holds at every scale with no error term, so it needs
no inflation, and a bound of the form $\Xi\le cY$ that it gives has $c(z)=c$.
In a bound of the second kind each structured energy is at most
$1+z$ times its main bound at a $z$-inflated counterexample; the
coefficients are nonnegative, so the sum of the main bounds is multiplied by
at most $1+z$; and the remainders, at most $zY$ in total, enter with
coefficients at most $c'$, which by \cref{def:robust}(a) is independent of
$\Lsc$ and $A$. If $\Xi\ge\Xi_0-U$ with $\Xi_0$ exact and $U\le u(z)Y$, then
$\Xi\ge\Xi_0-u(z)Y$. In the last case, by continuity there is $z_0>0$ such
that the strict inequalities hold for $0\le z\le z_0$; the argument then
applies at every $z$-inflated counterexample with $z\le z_0$, and
$F(c_1(z),\dots,c_n(z))\to F(c_1(0),\dots,c_n(0))$ as $z\to0$. The number of
remainders and their coefficients remain bounded, being those of finitely
many arguments.

(ii) We go through the proofs of these statements. The following are
\emph{exact relations}: $Y[\bigotimes_{\Blk}\mu_{\Blk}]=Y$; the orders
\ref{it:record:E2} and \ref{it:record:E3} and the relations
\ref{it:record:P1}--\ref{it:record:P5} of \cref{lem:energy-toolkit}, among
them the product centred interpolation \ref{it:record:P3b}; the
cut-order and bracket
manipulations (\cref{eq:cc:kron}, \cref{lem:record:block-gram}, part (i) of
\cref{lem:record:variance-interface}); the vanishing of the terms with an odd
number of skew factors in \cref{lem:record:record-expansion}; Step~1 of the
first comparison, which is (CP2)(a); the identity
\eqref{eq:record:expansion}; and the inequalities of
\cref{lem:record:common-mean} and \cref{cor:index-support-frame}, which hold
for arbitrary vectors and envelopes. The following are upper bounds of the second kind of (i):
\ref{it:record:E1}; the tagged expansion \cref{lem:sharp-first}, which bounds
each term of the tagged series, indexed by the numbers of
summands of the two tags, by the product of the sharp norms of
one tag times the geometric mean of two folded energies
(\cref{lem:record:folded-energy}(a),
(d$'$), (e)), the tag used for a category being chosen from the prices alone,
so that the bound is again a sum of structured energies with nonnegative
coefficients, and is used for Step~2 of \cref{lem:record:first-stage};
\cref{lem:bracket-energy}, the same expansion
with the bracket outside, in which the geometric mean of two inflated factors
of a Cauchy--Schwarz step is again inflated by $1+z$, so that $B$ becomes
$(1+z)B$ and not $(1+z)^2B$; the small-union bound underlying
$|I_\omega|\le(B+o(1))\pi_\omega Y$ of
\cref{lem:record:record-expansion}(i); the energy bounds of
\cref{lem:envelope-energy} and of the wrapping terms of
\cref{lem:record:wrap-expansion}, in which the conductors are held fixed
through every fold and averaged only at the end ((O3) of
\cref{def:admissible}); and the upper bound in
\cref{lem:record:variance-interface}(ii). Each of these is a sum of structured
energies with nonnegative coefficients bounded by convergent price products
such as $e^P\tagser(d,v)$, $e^{U_0}\tagser(d,v)$, $e^{S+A}$ and $e^T$. The
following are lower
bounds of the third kind: the first comparison $Y_\lambda\ge(1-L_0-o(1))Y$,
which is Step~1 minus Step~2; the bound $Y^{\Nbr}[J_X]\ge
Y[J_X]-(R_{\rm err}+o(1))C_XY$ of \cref{lem:record:variance-interface}(ii),
which is \ref{it:record:E2} minus \ref{it:record:E1}; and the gain
$\max(0,\ell/\bar C-R_{\rm err})$ of \cref{lem:record:variance-interface}(iii),
a continuous function of robust bounds. Finally,
\cref{lem:record:record-expansion}(ii), (iii) combine the bounds of (i) of
that lemma with the class bounds of (CP5), which are robust uniformly in the
class by (CP8)(b), through the complete polynomial $\Psi$. The series
$\sum_\omega b_\omega$ that $\Psi$ dominates has as coefficients $B$, the scalars $c_{\mathcal C}$ and the coefficients of the class bounds; at a $z$-inflated
counterexample these become functions of $z$ that converge to them uniformly
over the classes, and we may replace each by its positive part, while the
errors $o(1)\pi_\omega$ of (CP5) become $\epsilon(z)\pi_\omega$ with one function
$\epsilon(z)\to0$. Adding to $\Psi$ the positive parts of the changes of the
$b_\omega$, and $\epsilon(z)(e^{S(x)}\cosh A(x)-1)$, gives a series $\Psi_z$
with nonnegative coefficients that dominates the inflated bounds. The
difference $\Psi_z(\hat\pi,\hat y)-\Psi(\hat\pi,\hat y)$ is at most the
largest change of a coefficient times the sum of the values at $(\hat\pi,\hat
y)$ of the monomials of the $b_\omega$, plus $\epsilon(z)(e^{\hat S}\cosh\hat
A-1)$. That sum is finite: the monomials in the prices sum to at most
$e^{\hat S}\cosh\hat A-1$, since distinct selections have distinct squarefree
monomials, and those of the class bounds to at most
$\prod_\Blk(1+\sum_j\hat y_{\Blk,j})$, since each of them contains one auxiliary
price of each atom of its selection and determines the selection (CP5); this
product is finite because $\hat y$ has a finite total. So $\Psi_z(\hat\pi,\hat
y)$ tends to $\Psi(\hat\pi,\hat y)$ as $z\to0$.

The remainders. By \cref{lem:record:wrap-expansion}, whose proof splits the
term of an index subset $\mathcal S$ by the set of its wrapping
queries (each query is of exactly one of three kinds: in
$\mathcal S$ and wrapping, in $\mathcal S$ and nonwrapping, or
outside $\mathcal S$), the nonwrapping terms are at most
$3^kn_\nu$ integrals of the
form bounded in \cref{lem:record:density-can}, one for each product component
of $\nu$, against the conditional law of the indices, with coefficients
adding up to at most $2^k\mathfrak v_\nu e^{P+S}$ for each pair; that proof
also bounds the mass of $\nu$ by (CP2) and the brackets of (CP4). This is
(a). The
selections with $q_E>N^\tau$ occur in finitely many expansions, each of the
form bounded in \cref{lem:record:density-lp} with post-reduction prices of
the selected laws as the $\eta_{\Blk}$. Each law chosen at an atom in these expansions is one
of the laws of (CP2)--(CP6) or a fold or normalization of one, so its price
is at most the sum of the prices of the atom listed in (CP8)(c); the
components of a signed comparison law are never selected laws, since they
occur only in the cancellation of (a), so the variation factor $\mathfrak
v_\nu$ does not enter here. Taking that
sum as $\eta_{\Blk}$ for all the expansions,
$\Pi_{\rm H}=\prod_{\Blk}(1+\eta_{\Blk}q_{\Blk}^{1/4})$ is finite by (CP8)(c);
summing the finitely many sums gives the bound of one such sum. The number of
terms and their coefficients do not depend on $\Lsc$ or $A$. This is (b).

(iii) Let $(\Lsc,A)$ be a $z$-inflated least normalized counterexample, with
$\Lsc$ and $z$ in the ranges of the finitely many robust bounds used. By (ii),
the steps of the proof of \cref{thm:record:criterion} hold with each $o(1)$
replaced by the deviation of a function continuous at $0$:
$2^kY_\nu\le J_\nu+2^kW(z)Y$; \eqref{eq:cp:mean-step} exactly; $V_{\rm L},V_{\rm
R}\le2^kJ(z)Y$; and $Y_\lambda\ge(1-L_0(z))Y$, with $W(0)=W$, $J(0)=J$ and
$L_0(0)=L_0$. By hypothesis, $Y_\nu\ge Y_\lambda-\mathcal E(z)Y$ and
$\frac12(e_{\rm L}+e_{\rm R})\ge m(z)Y$, with $\mathcal E(0)=\mathcal E$ and
$m(0)=m$. The algebra of that proof gives $Y\le\Theta(z)Y$ with
\[
 \Theta(z)=\bigl(1-2^{-k}\bigr)\bigl(L_0(z)+\mathcal E(z)\bigr)+J(z)+W(z)
 +\frac{1-m(z)}{2^k},
\]
which is continuous at $0$ with $\Theta(0)=\Theta(\mathcal E,m)$; since
$Y>0$, the inequality persists when $\Theta(z)$ is replaced by
$\max(\Theta(z),0)$. This is (Rec), with $z_0$ the minimum and $\Lsc_1$ the maximum of the numbers $z_0,\Lsc_1$ of the robust bounds involved. Its
remainders are those of (ii) and those of the robust arguments for (i) and
(ii) of \cref{thm:record:criterion}; by \cref{def:robust}(a) they are
finitely many integrals of the form of \cref{lem:record:density-can}, with
coefficients bounded independently of $\Lsc$ and $A$, and finitely many sums
of the form of \cref{lem:record:density-lp}, which together have the bound of
one such sum with the sum of their constants $c$ and the largest of their
products $\Pi_{\rm H}$. This is the second hypothesis of
\cref{thm:record:lineardensity}.
\end{proof}

\begin{remark}[Robustness of the final chain]\label{rem:cp:robust-final}
In Part~V every bound that enters (CP5) and (CP7) is of one of the kinds of
\cref{lem:record:rec}(i).
\begin{enumerate}[label=(\alph*)]
\item\emph{Upper bounds of structured energies}, with nonnegative
 coefficients over certified prices: the direct prices of the middle records
 and their repricing (\cref{lem:middle-comparison-prices,lem:record:repricing});
 the repriced first-stage totals and the tagged series, which give $U_0$,
 $B$, $R_{\rm err}$ and $L_0$ (\cref{cor:first-stage-repricing,cor:two-cut-series});
 every leading class cost (with $n,m$ and $\kappa_H,\kappa_b,\Uinv,\delta_{\rm s}$ as in \cref{cor:envelope-costs})
 \[
  C_{n,m}=\kappa_He^{U_0}-(\kappa_H-\kappa_b)e^{t\Uinv}+M_{n,m}R_{\rm err}
  -M^{\rm inv}_{n,m}e^{t\Uinv}t\delta_{\rm s}
 \]
 whose terms are an outside energy, the
 invariant-background polynomial with nonnegative coefficients
 (\cref{lem:oct24:invariant-background,rem:completion-subtraction}),
 normalized couplings bounded by \ref{it:record:E1}, and an orbit coupon that
 lowers the coupling term by an exact orbit decomposition
 (\cref{lem:orbit-coupon}); the remainder polynomials $B\mathcal R_{n,m}$; $J$
 and $W$; the upper prices $S_{\rm u}(a)$, $A_{\rm u}(a)$
 (\cref{prop:record:upper}); and the complete polynomial after the
 repeated-block monomials are added with nonnegative coefficients
 (\cref{lem:record:record-expansion}(iv), (v)).
\item\emph{Exact same-scale relations}: the orders \ref{it:record:E2},
 \ref{it:record:E3} and \ref{it:record:P3b}; the operator comparison of the
 pair projections with their positive substitutes
 (\cref{lem:oct24:pure-pair-substitute}); the matrix-envelope orders and their
 mask form (\cref{lem:oct24:mask-schur}); the changes of variables of
 \cref{lem:reflection-averaging}; the orbit decomposition of
 \cref{lem:orbit-coupon}; the positive compressions of the aligned matching
 bridge (\cref{lem:record:paired-e-joint}(a), (c)); the all-degree
 polarization identity (\cref{lem:tr:polarization}); and the contraction and
 density steps of \cref{lem:record:companion-energy}.
\item\emph{Exact relations minus upper bounds}: the restorations $l$ with the
 individual piece caps (\cref{lem:record:restoration,lem:piece-caps}; under
 inflation $l=\max(0,1-(C-1)h)$ becomes $\max(0,1-((1+z)C-1)h)$); the gains
 with positive backgrounds and the late-pool gains
 (\cref{prop:background-variance,lem:pool-coupons,lem:var-gain}), positive
 parts of continuous expressions; the selected-only aligned matching gain
 (\cref{lem:record:paired-e-joint,lem:class-credit}); the credited class
 bounds (\cref{lem:class-credit}); the variance credit $D_{\rm var}$
 (\cref{cor:variance:final}); and the mean credit $m$
 (\cref{prop:record:mean-credit}), which subtracts $R_{\rm
 err}e^{T_K+T_{QE}}$.
\end{enumerate}
The strict inequalities used are: every saving used in (a)--(c) has a stated positive lower bound; every restoration gain lies in $[0,1)$ with a positive denominator; the
two-cut conditions $d<6$, $dv<36$ of \cref{cor:two-cut-series}; and the
positive-mean branch; the closing and restoration scripts assert them with
strict margins (\cref{tab:cert:finite}). Three further gates are
nonnegativity conditions that need no margin. The orbit-coupon gate
$M^{\rm inv}_{n,m}e^{t\Uinv}t\delta_{\rm s}\le M_{n,m}R_{\rm err}$ concerns the
coefficient $M_{n,m}R_{\rm err}-M^{\rm inv}_{n,m}e^{t\Uinv}t\delta_{\rm s}$ of
the coupling term, which is itself a bound of the kind (a) and so only grows
under inflation. Every subtracted term $cg$ on a monomial in the auxiliary prices $y^{\rm arm}$ of a replaced class,
a fixed envelope cap $c$ times a gain $g<1$, is below the coefficient
$C_{n,m}$ of that monomial (\cref{lem:class-credit}); this gate is proved in
the text, from $c\le\kappa_be^{t\Uinv}\le C_{n,m}$ (\cref{prop:closing-data}),
not asserted by a script. Under inflation the coefficient, a bound of the kind
(a), only grows, and the credit, its gain being of the kind (c), only shrinks.
The restoration regime $0\le l\le C/2$ is kept by inflation, which lowers
$l=\max\bigl(0,1-(C-1)h\bigr)$ and raises $C$. The side coefficients
$k_s\ge0$, the side increments $\delta^+_{s,X}$ of \eqref{eq:signed-increment} and the bracket condition (\cref{prop:signed-sides,prop:record:matching,eq:protection})
are inequalities between fixed numbers that enter only exact bracket
inequalities, so inflation does not affect them; the piece caps
$1\le C_i\le e^{2aL}$ (\cref{lem:piece-caps}) and the domain $0<a\le9/10000$
of the analytic gates are fixed numbers as well. Under inflation the class
costs and $B$ increase together and every normalized restoration stays at
most one. The selections of the classes and of the credits have at most six
atoms, the envelope expansions have finitely many couplings, the comparison
law has finitely many product components, and $\Kint$ is fixed, so the
remainders are finitely many.
\end{remark}

\subsection{End of the proof: applying the baseline induction}\label{subsec:cp:proof}

Everything is now in place: the record expansion bounds
$|Y_\nu-Y_\lambda|$, hypothesis~(CP7) supplies the mean credit $m$,
and the closing lemma together with the induction of
\cref{sec:scale-comparison} finishes Theorem~C.

\begin{proof}[Proof of \cref{thm:comparison}]
Let $(\Lsc,A)$ be a least normalized counterexample. By
\cref{lem:record:record-expansion}(ii), applied with the classes and class
bounds of (CP5), the upper prices $\hat\pi$ and the upper auxiliary prices
$\hat y$,
\[
 |Y_\nu-Y_\lambda|\le(\Psi(\hat\pi,\hat y)+o(1))Y\le(\mathcal E+o(1))Y,
\]
which is hypothesis (i) of \cref{thm:record:criterion}; (CP7) is hypothesis
(ii). Both hold robustly: the first by \cref{lem:record:rec}(ii), with
$\mathcal E(z)=\mathcal E+|\Psi_z(\hat\pi,\hat y)-\Psi(\hat\pi,\hat y)|+c'z$,
where $\Psi_z$ is the inflated series of that proof and $c'$ bounds the
coefficients of the remainders (\cref{lem:record:rec}(i)), and the second by
(CP8)(b). By \cref{lem:record:rec}(iii) there are $z_0>0$, $\Lsc_1$ and a
function $\Theta(\cdot)$, continuous at $0$ with
$\Theta(0)=\max(\Theta,0)<1$, for which (Rec) of
\cref{thm:record:lineardensity} holds, with remainders of the forms that
theorem requires. Its other hypotheses hold: $0<a<\sigma$ and
$0<\theta\le1$ by (CP0)(b), and \ref{it:record:R1}--\ref{it:record:R5} by
(CP1) and (CP4)(a), with the clauses of \ref{it:record:R4} and
\ref{it:record:R5} on the accepted parts read for the normalized components
of $\nu$ and the queries of $\mathcal E_k$, as (CP4)(a) states them; these
clauses enter only through the hypotheses of \cref{lem:record:density-can},
which \cref{lem:record:wrap-expansion} checks for every product
component. It gives a constant $C$, depending only on the
comparison data, with
$\Yf(\Lsc,A)\le C\Lsc^{-a}\alpha^\theta$ for every $\Lsc\ge1$ and every
square-difference-free $A\subseteq[\Kint\Lsc]$, and
$s(N)\le C'N^{1-a/(d-\theta)}$ for every $N\ge1$, with $d=72$.
\end{proof}

The proof is the composition of \cref{thm:record:criterion},
\cref{lem:record:first-stage}, \cref{lem:bracket-energy} (inside
\cref{lem:record:record-expansion,lem:envelope-energy}),
\cref{lem:record:wrap-expansion}, \cref{lem:record:record-expansion},
\cref{lem:record:common-mean} (through \cref{cor:index-support-frame}),
\cref{lem:envelope-energy} and \cref{lem:record:rec}, followed by
\cref{thm:record:lineardensity} of Part~II, and of nothing else. No symmetry
of $\nu$ and no positivity of $\nu$ is used, except that each of its finitely
many product components is a nonnegative accepting law in the cancellation of
the nonwrapping terms (\cref{lem:record:wrap-expansion}); the skew parts enter
only through the vanishing of the terms with an odd number of skew factors in the record expansion and through
the separate pricing of skew changes.

\subsection{The case of one query (\texorpdfstring{\cref{prop:mech:one-query}}{Proposition \ref*{prop:mech:one-query}})}\label{subsec:cp:one-query}

For completeness we record how the one-query proposition of
Part~I follows from the same argument with $k=1$. Nothing later
depends on this remark.

\begin{remark}[Proof of the one-query proposition]\label{rem:cp:one-query}
We read the hypotheses of \cref{prop:mech:one-query} with the notation of this part: $k=1$ and $\mathcal E_1=\{e\}$; $\mu$ satisfies (CP1); $\nu$ is a
product of nonnegative atom measures of bounded total mass satisfying
(CP4)(a) for the query $e$; (H1) is
$Y_\nu\ge(1-\mathcal E-o(1))Y$; and (H2) gives nonnegative measures
$M_{\Blk}$ with $\bigl(\begin{smallmatrix}M_{\Blk}&\nu_{\Blk}\\
\nu_{\Blk}^*&M_{\Blk}\end{smallmatrix}\bigr)\succeq_\iota0$ such that the
energy under $\bigotimes_{\Blk}M_{\Blk}$ of every structured slot datum
(\cref{def:record:structured-energy}(c)) is at most $(B+o(1))$ times its
coarse price times $Y$. With $k=1$, \eqref{eq:record:expansion} reads
$2Y_\nu=J_\nu+I_{\nu,\{e\}}$.

(1) Part (b) of \cref{lem:record:common-mean}, applied as in
\cref{cor:index-support-frame} with $M_{\rm L}=M_{\rm R}$ the cut matrix of
$\bigotimes M_{\Blk}$, $\Gamma$ that of $\nu$, and $\mathcal T=2P_e-T_e$, gives
$|J_\nu|\le V_{\rm L}^{1/2}V_{\rm R}^{1/2}$ with $V_{\rm L}=V_{\bigotimes M}(X)$ and
$V_{\rm R}=V_{\bigotimes M}(Y)$. By \cref{lem:twointerval} the energy of
$X^{\Jc}$ under $\bigotimes M_{\Blk}$ is the energy under this measure of a
structured slot datum whose changed slots are $u_e$ and $\iota u_e$, each
with the slot price
$(\Kint\sq(D_{\Jc_e}))^{a/d}$ and the coarse factor $1$ of a uniform index; by
(H2) it is at most $(B+o(1))(\Kint\sq(D_{\Jc_e}))^{2a/d}Y$, the conductor
being held fixed and averaged at the end. Averaging over $\Jc$ gives $V_{\rm
L}\le(BR_I+o(1))Y$, and likewise for $V_{\rm R}$; so $|J_\nu|\le(BR_I+o(1))Y$.

(2) Split $I_{\nu,\{e\}}$ according to $b_e$. If $b_e\ne\Kint$, the query $e$
is accepted, its two ends carry $g^{I_{b_e},\Mzero,s_e}$ and
$g^{I_{b_e+1},\Mzero,s_e+r_e}$ with ordered supports and a nonvanishing square class
$r_e$ (\cref{lem:imp:padic-envelope,def:nonvanishing-square-class}), and the term is a cancellation remainder of
the form of \cref{lem:record:density-can}, hence $o(1)Y$. The event
$b_e=\Kint$ has probability $1/\Kint$. At this fixed interval,
the form of \cref{lem:record:common-mean} for the filtration by
residue coordinates, with the operator
$T_e$, bounds the term by the geometric mean of two energies as in (1), in
which the fixed interval adds the coarse factor $\Kint^{\theta/d}$ at each of
the two changed slots. Hence $|I_{\nu,\{e\}}|\le(BR_I\Kint^{2\theta/d}/\Kint
+o(1))Y=(2W_1+o(1))Y$ with $W_1=BR_I/(2\Kint^{1-2\theta/d})$.

(3) Therefore $2(1-\mathcal E-o(1))Y\le2Y_\nu=J_\nu+I_{\nu,\{e\}}\le
(BR_I+2W_1+o(1))Y$, that is, $1\le\mathcal E+BR_I/2+W_1+o(1)$, which is the
assertion of \cref{prop:mech:one-query}, where $\theta=1$ and
$\Kint^{1-2\theta/d}=\Kint^{35/36}$.

If part (a) of \cref{lem:record:common-mean} is used in (1) instead, then
$J_\nu\le Y_\nu+(BR_I+o(1))Y-\frac12(e_{\rm L}+e_{\rm R})$, and with
$\frac12(e_{\rm L}+e_{\rm R})\ge(m-o(1))Y$ one obtains $1\le\frac12\mathcal
E+BR_I/2+W_1+(1-m)/2+o(1)$: keeping the common mean halves the coefficient of
$\mathcal E$. This is the case $k=1$ of \cref{thm:comparison}, without a first
stage ($L_0=0$) and with the one energy factor $B$ of (H2) in place of both
$Be^T$ in $J$ and $Be^{S+A}$ in $W$.
\end{remark}

\subsection{Where the hypotheses are verified (table for Theorem~C)}\label{subsec:cp:table}

\Cref{tab:hypotheses} lists each hypothesis of Theorem~C, what Part~V shows
for the construction of \cref{def:record:parameters}, and the statements that
show it. Each section of Part~V ends with a paragraph naming the rows it
verifies.

{\small
\begin{longtable}{@{}>{\raggedright\arraybackslash}p{2.75cm}>{\raggedright\arraybackslash}p{6.1cm}>{\raggedright\arraybackslash}p{5.95cm}@{}}
\caption{Where the hypotheses of Theorem~C are verified.}\label{tab:hypotheses}\\
\toprule
Hypothesis & What is shown & Proved in\\
\midrule
\endfirsthead
\toprule
Hypothesis & What is shown & Proved in\\
\midrule
\endhead
\bottomrule
\endfoot
(CP0)(a) slots, cuts & $V$ sorted by seven folds; $\iota$ fixed-point free; the family $\mathfrak F$ & \cref{def:slot-array,lem:slot-array}\\
(CP0)(b) atoms, exponents, $a<\sigma$ & $a=\PushA<9/10000<\sigma=1/825$; inequalities (S1)--(S7); $\Kint=\PushK$; checked inequalities involving $a$ hold for $0<a\le9/10000$ & \cref{def:record:parameters,def:record:functional,cert:record:analytic-stage,cert:record:frame-data}; data files in \cref{app:certificates}\\
(CP0)(c) queries & six queries with disjoint endpoints, sources and reflected targets in $V_+$ & \cref{lem:slot-array}, \cref{subsec:front6:construction}\\
(CP0)(d) index law & finite conductor depths $m_{\rm ext}=64$; nonvanishing square increments modulo $\Mzero$; $R_I=M_{\rm pad}(\Kint)<1.247139907407$ & \cref{lem:kernel-depth-classes,lem:imp:padic-envelope,lem:hyb:moment,prop:imp:padic-price,cor:imp:padic-majorant,cert:index-moment}\\
(CP1) prefix & probability law, reflection positive at every cut, square-affine invariant & \cref{lem:frontier:iidlocal}\\
(CP1) middle & the same for the twelve-label reference; the uniform root below $2^{30}$, and seven positive roots from $2^{30}$ & \cref{lem:front12:reference} via \cref{lem:five-type}; root caps by \cref{lem:layer12-caps} (route A, \cref{thm:spectral-criterion}), \cref{prop:l20-collisions,prop:root-chamber,prop:root-layer}; the fold cap \cref{lem:old-f-law}; rows by \cref{prop:middle-negative,lem:row-comparison,lem:four-count}\\
(CP1) upper & the same for the two-layer reference in $310$ bands & \cref{lem:upper-reference-rp}, \cref{lem:five-type} (with (f)); \cref{thm:spectral-criterion,thm:record:direct-star,lem:logconvex-coefficients,prop:observed-caps,lem:joint-rows,lem:joint-pins}\\
(CP1) tail & the same for the large-prime law, at primes of both classes mod $4$ & \cref{lem:tail-law} (the tail law of \cref{def:tail-law}); class-one caps, \cref{subsec:advance:parity}\\
(R4) & $(\gamma,v)$-mixing, $\gamma=255/1024$, at $p\ge P_5$; query pairs & \cref{lem:record:tail-mixing,lem:trace-mixing,lem:mixture}; (S7) of \cref{cert:record:analytic-stage}\\
(R5), (CP4)(a) & accepted query pairs are strict-square pairs for every component; clique differences of even valuation with square unit part; uniform one-slot marginals & \cref{lem:frontier:iidlocal,cert:record:middle-geometry,lem:signed-components,def:record:mixture,lem:record:e-law,lem:merge-repair,cor:merge-repair}, \cref{sec:frontier:upper-oneway}, \cref{lem:six-query-retention,def:dense-comparison,prop:le-mixture}, \cref{lem:tail-prices}(i)\\
(CP2)(a) & $E_{\rm add}+G-F_{\rm fold}\succeq_{\vartheta_0}0$ & \cref{lem:middle-polarization}: \cref{lem:tr:polarization} with the rows of \cref{lem:front12:reference}(i) (\cref{prop:middle-negative,lem:row-comparison})\\
(CP2)(b) & $E_{\rm new}$ keeps the actual conditional marginals; two tags, $E_0$ sharp at its first cut and $E_1$ at its own, with equal sharp and equal variation prices; sharp coefficient $k_{\rm new}$ & \cref{lem:record:one-stage-folded-e,cert:record:one-stage}\\
(CP3) & $\Nbr\pm\lambda\succeq_\iota0$ for $\Nbr=\lambda+F_{\rm fold}+H$; $\Nbrz-\mu=E_{\rm add}+G+H$ positive reserved diagonal; the protection gate $v_D\ge2w_Dk_D$ & \cref{lem:record:gram}(i), (ii); \cref{lem:front12:reference}(iii), and (iv) for the protection gate (\cref{eq:protection}); \cref{sec:frontier:empty-triple-negative}(c)\\
$P$, $V$, $D$, $U_0$, tags & the middle-frame prices $P,V,D$, the block prices, repricing, the tagged series & \cref{lem:record:middle-prices,lem:record:block-charge,cert:record:middle-price,lem:record:repricing,cor:first-stage-repricing,cor:two-cut-series}; \cref{prop:closing-data}(a)\\
(CP4) middle $F$ & the signed whole-block law: three positive accepting components repaired to the actual marginal, exact marginals, variation $1+\tlaw$; direct prices & \cref{lem:middle-comparison-prices} with \cref{def:record:mixture,lem:signed-components,prop:signed-bounds,lem:actual-features,cor:markov-completion}\\
(CP4) middle $E$ & the law $K_1$ with exact marginals; prices; degree-four remainder & \cref{lem:middle-comparison-prices} with \cref{lem:record:e-law,lem:record:exact-marginal-e,lem:oct24:matching-star}\\
(CP4) merges, deletions & asymmetric repairs; deleted positive records & \cref{lem:middle-comparison-prices} with \cref{lem:merge-repair,cor:merge-repair,rem:cp:charging}\\
(CP4) upper & $84$ one-way repairs (skew-only charging; forest records in the bands from $2^{48}$), comparisons by the dense EL law, the LE mixture, deletions; $S_{\rm u}(a)$, $A_{\rm u}(a)$ & \cref{prop:record:upper,lem:tr:repair,lem:six-query-retention,prop:dense-comparison,prop:le-mixture,prop:forest-records}, \cref{rem:cp:charging}(iii); \cref{lem:upper-stars} (tail)\\
(CP4) prefix & the deleted exception, reserved diagonal, charged to $S$ & \cref{lem:frontier:iidenergy}\\
(CP4) tail & repairs with the all-prime caps and deletions; $s_p+a_p\le C_{\rm tail}p^{-3/2+2a}$; $S_{\rm tail}$, $A_{\rm tail}$ at $a$ & \cref{lem:tail-prices}\\
(CP5) variance classes & the variance savings with positive backgrounds, of degrees one to four, and the savings pooled over the middle intervals with least prime at least $2^{31}$, with the piece caps & \cref{cor:variance:final} via \cref{lem:var-gain,prop:background-variance,lem:pool-coupons,lem:record:restoration,lem:piece-caps,lem:record:variance-interface,lem:record:paired-e-joint}\\
(CP5) envelope classes & the class bounds $C_{n,m}(A^{\rm arm}_F)^n(A^{\rm arm}_E)^m+B\mathcal R_{n,m}$ ($n\ge1$, $n+m\in\{2,4,6\}$), of the aligned and of the crossed matching pairs, and of $(0,4)$, $(0,6)$ & \cref{cor:envelope-costs} via \cref{thm:envelopes,lem:oct24:mask-schur,lem:reflection-averaging,lem:oct24:invariant-background,lem:orbit-coupon}; \cref{lem:record:record-expansion}(v)\\
(CP5) credited classes & $\kappa_{FF}g$ on fold pairs; the selected-only $\kappa_bg$ on aligned matching pairs; the quartic alternatives & \cref{lem:class-credit}\\
(CP6) & the matching bracket with the signed direct Gram, the sides $\max(\bar H_Xd_\tlaw,d^{\rm aux}_s)$ and the chosen diagonals; $T$ & \cref{prop:record:matching,prop:signed-sides,lem:record:direct-gram,lem:record:asymmetric-gram}\\
(CP7) & the mean credit $m=1+m_F+m_Q-R_{\rm err}e^{T_K+T_{QE}}$ & \cref{prop:record:mean-credit,lem:record:companion-energy}\\
(CP8)(a) & (O1)--(O4) for every selected law & \cref{def:admissible} and \cref{tab:admissibility}\\
(CP8)(b) & every bound is of a kind of \cref{lem:record:rec}(i); strict gates & \cref{lem:record:rec,rem:cp:robust-final}; \cref{thm:envelopes}(iv)\\
(CP8)(c) & at a tail atom only $s_p$, $a_p$ and $T_{{\rm L},p}=T_{{\rm R},p}$ are nonzero, with $T_{{\rm L},p}\le s_p+a_p$ ($M_{s,p}-\mu_p$ is the sum of the norm envelopes), and $\sum_p(s_p+a_p)p^{1/4}<\infty$; auxiliary prices only at middle atoms; the other atoms are finitely many & \cref{lem:tail-prices}(iii), (iv); \cref{prop:record:matching}(i)\\
$\Theta<1$ (baseline) & $\Theta<\PushTheta$ & \cref{prop:closing-data,cert:closing}; data files in \cref{app:certificates}\\
\end{longtable}
}

\paragraph{Consequences and what to read next.}
The lemmas of the previous section and \ref{it:record:E1}--\ref{it:record:E3};
(CP0)--(CP8) with \cref{def:hyb:coupled,lem:func:unitsq,def:robust}
and \eqref{eq:record:expansion}; and from Part~II
\cref{prop:record:density-structured,lem:record:folded-energy,lem:record:density-can,lem:record:density-lp,thm:record:lineardensity},
\cref{def:admissible,lem:record:block-gram} and \cref{eq:cc:kron}.
